\section{Results}
\label{sec:coral-island-results}

The resulting model for coral reef island generation enables a high level of user control as unconstrained painting allows complex scenarios while producing height fields in real-time. In this paper we used Blender to provide renders directly from the output height fields. As our pix2pix model is trained to output $256\times256$ images, the resolution of the 3D models is limited by this architecture.

\begin{figure}
    \autofitgraphics[]{3_features.png, 3_heightmaps.png, 3_results.png}
    \autofitgraphics[]{4_features.png, 4_heightmaps.png, 4_results.png}
    \autofitcaptions{Input sketches, Output height field, Result}
    \caption[Output of the cGAN on fuzzy inputs]{User applied fuzzy brushes to draw the label map, resulting in some pixels that are inconsistent with the dataset and illogical island layouts: abyss regions (red) are found between beach (green), lagoon (cyan) and reef regions (blue). The neural model ignores the layout inconsistencies, and partial over-brightness as long as the pixel is not completely white.}
    \label{fig:coral-island-results-fuzzy}
\end{figure}
\begin{figure}
    \autofitgraphics[]{DinoIsland_features.png, DinoIsland_heightmaps.png, DinoIsland_results.png}
    \autofitcaptions{Input sketches, Output height field, Result}
    \caption[Output of the cGAN with new colors in the input]{Without constraints on the generation, the user may use unrealistic layout and the neural network will however output a plausible result. Here new colours have been added (pistachio), but the network naturally considers it as a transition from mountain to beach.}
    \label{fig:coral-island-results-dino}
\end{figure}


By using deep-learning-based models, most initial constraints are removed (radial layout, isolated islands, ...). The control over the overall shapes of the island regions is given through digital painting with any software, here using GIMP (first column of \cref{fig:coral-island-results-fuzzy}) and Paint (such as in \cref{fig:coral-island-results-dino}). Each pixel of the image is encoded in HSV, with the region identifier encoded in the Hue channel. The user may increase or decrease the subsidence level of the island by modifying the Value channel over the whole image (see \cref{fig:coral-island-results-subsidence}).

Since the model relies on pixel-level statistics rather than strict class values, it is robust to noise caused by compression, anti-aliasing from brush tools, or resizing artefacts introduced by image editors. The example displayed in \cref{fig:coral-island-results-fuzzy} (top) displays a sketch with blurry outlines and unexpected layouts (see the small red regions inside the southern lagoon region or the adjacency of beach regions directly with the abyssal region). The learned model does not include inconsistencies and results in plausible 3D models. We can note that the input label map is not required to be hand-drawn, as a simple Perlin noise can produce it randomly, as shown in the examples of \cref{fig:coral-island-perlin-examples}. \cref{fig:coral-island-results-fuzzy} (bottom) shows highlight clipping (white pixels due to artificial oversaturation, losing the Hue value) that is not interpretable by the model and resulting in black patches in the result. 

The tolerance to imprecise input values enables more control about the transitions between two regions than the procedural module. \cref{fig:coral-island-results-dino} shows an example of input map with regions that are leaking over neighbouring regions, and the introduction of new hue values nonexistent in the dataset (light green and dark green) but are the interpolated hue values of mountain regions and beach regions.

Because our curve-based procedural phase included low randomness, the output of the cGAN is limiting its unpredictability, meaning that small input changes result only in minor changes on the output, preventing unexpected results. \cref{fig:coral-island-results-subsidence} shows the result of an input map with only a variation on the subsidence level, the resulting height fields are very similar. Adding the real-time computation of outputs, it becomes possible to construct progressively a landscape and correct small inconsistencies to intuitively design islands inspired by real-world regions. A creation of an island similar to Mayotte, presented in \cref{fig:coral-island-example-Mayotte} was hand-made in just a few minutes. 
% The islands presented at the beginning of the chapter in \cref{fig:coral-island-island-examples} are also reproduced in \cref{fig:coral-island-results-reproductions} with different subsidence rates.

\cref{fig:coral-island-results-volcanic-training} (bottom) presents results of a second cGAN training for which the procedural parameters used for the generation of the training dataset "volcanic" are slightly adjusted for lower overall resistance and to have a crater in the center of islands (\cref{fig:coral-island-results-volcanic-training}, top). The $\heightProfile$ function include a hole at the center of the island to emulate a volcaninc crater, which we can see strongly in the first and last examples, and more slightly in the second example; the resistance function $\resistance$ is lowered on the reefs, creating visible dents in the reef and dendritic patterns on the island sides.

\begin{figure}
    \autofitgraphics[]{Mayotte-example.png, 5_features.png, 5_results.png}
    \caption[Comparison of  the cGAN output with Mayotte Island]{Comparison between real (left, from Google Earth) and synthetic (right) islands of Mayotte. The label map (centre) is hand-drawn to match approximately its real-world counterpart.}
    \label{fig:coral-island-example-Mayotte}
\end{figure}




\subsection{Advantages of the approach}
\label{sec:coral-island-advantages}

One of the main strengths of this approach is its ability to produce a wide variety of island terrains, even in the absence of real-world data. The procedural generation methods allow for high flexibility in designing both the shape and features of the island, while the use of cGAN enables further refinement and the generation of terrains not bound by the original constraints of the procedural model. By combining these two methods, we leverage the complementary advantages of both: the structured control of procedural techniques and the pattern-learning capabilities of deep learning.

A key advantage of this approach is the preservation of semantic information throughout the generation process. The label map, which serves as the input to the cGAN, can also be used after terrain generation to provide a detailed semantic representation of the different regions of the island (island, beach, lagoon, coral reef, and island body). This label map can be useful to guide post-processing operations, such as applying different textures based on terrain features or adding other environmental elements, vegetation for instance. The preservation of semantic information provides a useful connection to the next stage of terrain manipulation, making the process more versatile and adaptable to different use cases.

Furthermore, the use of an out-of-the-box cGAN model highlights the feasibility of employing existing neural network architectures with minimal modifications in the field of procedural generation. This is particularly important for domains where real-world data is scarce, such as coral reef islands, allowing synthetic data to be effectively used for training purposes.

\subsection{Limitations}
\label{sec:coral-island-limitations}

While the cGAN model provides increased flexibility and variety in island generation, it does come with certain limitations:

\begin{Itemize}
    \Item{} \textbf{Data-driven dependence:} Just as any deep learning method, the quality of generated terrains strongly depends on the quality of the dataset. Unusual layouts or out-of-distribution inputs are not entirely interpretable by the models. Our dataset generation method lifts some restrictions, but some remain such as the required presence of all regions in the input map. This limitation reduces the true diversity of the generated terrains, as the cGAN cannot generate terrains that deviate too far from the examples in the training set.
    % The quality of the generated terrain depends entirely on the quality of the training dataset. Since the dataset is synthetically generated, any limitations or biases in the initial dataset directly affect the cGAN's output. This dependence on data quality makes it crucial to design a well-augmented and varied dataset to ensure diverse and realistic outputs.
    \Item{} \textbf{Biases from the synthetic dataset:} Since the cGAN model is trained entirely on a procedurally generated dataset, it inherits the biases present in the initial algorithm. For example, while the model breaks free from the radial symmetry constraint and centre positioning, it remains dependant on the structure and patterns of the synthetic data (e.g., our procedural drainage patterns, our reef analytic morphology, etc.). %This limitation reduces the true diversity of the generated terrains, as the cGAN cannot generate terrains that deviate too far from the examples in the training set.
    The "realistic" aspect of the results thus depends on the generation methods used to create the dataset.
    \Item{} \textbf{Lack of user control:} Another limitation of using cGAN in this context is the lack of real-time user control during terrain generation. While traditional procedural generation methods allow users to adjust parameters during the generation process, the cGAN model operation is abstracted from the user, providing no mechanism for direct interaction beyond the initial label map. This reduces the level of customisation available to the user.
\end{Itemize}